% Teaching additions; keys refer to the private optics layout, not printed numbers.
%<append:棱镜偏向、色散与彩虹>
\Needspace{8\baselineskip}\begin{origexamplebox}{固定棱镜中的不同颜色}
\label{ch27:ex:depth-abbe}
光学玻璃的阿贝数定义为 $\nu=(n_D-1)/(n_F-n_C)$。顶角 $A=60^\circ$ 的棱镜有 $n_D=1.500$、$n_C=1.495$、$\nu=40$，置于空气中。白光以同一方向入射，调节棱镜，使 $D$ 光在棱镜内的传播方向垂直于顶角平分面。此后固定棱镜，求出射的 $C$ 光和 $F$ 光之间的夹角。
\end{origexamplebox}\par\nopagebreak[3]
\begin{xesolution}
先区分三种折射率。阿贝数给出的是两条谱线的折射率差，而不是折射率与波长的比值：
\[
 n_F-n_C=\frac{n_D-1}{\nu}=\frac{0.500}{40}=0.0125,
 \qquad n_F=1.5075.
\]
$D$ 光沿对称光路传播，因此在两个侧面处的内部折射角都是 $A/2=30^\circ$。参照图~\ref{ch27:fig:s42-19-6} 右图的两个工作面，分别相对该面法线计算折射角，第一面的共同入射角为
\[
 \sin i=n_D\sin(A/2)=0.750,
 \qquad i=48.5904^\circ.
\]
固定棱镜以后，$i$ 对三种颜色相同，但内部折射角并不相同。对任意折射率 $n$，逐次计算
\[
 r_1(n)=\arcsin\!\frac{\sin i}{n},\qquad
 r_2(n)=A-r_1(n),\qquad
 e(n)=\arcsin\!\bigl[n\sin r_2(n)\bigr].
\]
角 $e$ 从第二面外法线量起。两种颜色都从第二面透出，因为各自的 $n\sin r_2<1$；不能省略这一步存在性检查。代入得
\[
 e_C=48.1589^\circ,\qquad e_F=49.2438^\circ,
 \qquad\boxed{e_F-e_C=1.0849^\circ}.
\]
两束出射光共用同一个第二面法线，故出射角之差就是它们的夹角；总偏向角中的共同 $i-A$ 会抵消。$n_F>n_C$ 对应更大的偏向，与正常色散的方向一致。若令 $n_F=n_C$，两条计算链完全相同，夹角应为零。

本题 $A=60^\circ$，不能用薄棱镜式 $\delta\simeq(n-1)A$ 代替精确计算；也不能分别把 $C,F$ 两色都取为各自的最小偏向光路，因为实验中棱镜只调节了一次。
\end{xesolution}
%</append>
%<block:光线模型、波面与光学发展>
\subsection{从波长、折射率到光程}
光的频率记为 $\nu$，角频率为 $\omega=2\pi\nu$。真空波长 $\lambda_0=c/\nu$ 与介质内波长不是同一个量。在静止界面上，入射波、反射波和折射波必须以相同频率振动，否则界面上的电场边界条件不能在每个时刻同时满足。若透明介质的相折射率为 $n$，则相速度 $v_p=c/n$，介质内波长为
\[
 \lambda=\frac{v_p}{\nu}=\frac{\lambda_0}{n}.
\]
所以“进入玻璃后波长变短”不表示光的频率降低。色散指 $n$ 随频率而变；比较不同颜色时，要分别使用各自的折射率。

一段长度为 $s$ 的均匀介质，使单色波积累的传播相位为
\[
 ks=\frac{2\pi}{\lambda}s=\frac{2\pi}{\lambda_0}ns.
\]
因此把 $ns$ 称为这段路的光程：它把“走过多长”和“介质使相位变化多快”合在一起。依次通过若干介质时，光程相加为 $\sum_jn_js_j$；折射率连续变化时才把求和改为积分。光程的单位仍是米，折射率没有量纲。它不是把折射后的折线路径拉直所得到的几何长度。

本章的普通玻璃、水和空气先按各向同性介质处理。在这种介质中，波面法线、波矢与能量传播方向一致；晶体中的非寻常光则不一定如此，第30章再区分波法线与能流方向。脉冲通过明显色散的介质时，能量包络的传播时间涉及群速度，也不能直接把相折射率代入 $t=ns/c$。
%</block>
%<block:反射定律与镜面光路>
\subsection{先定法线，再追踪反射方向}
入射角从法线量起，不从镜面量起。例如光线与镜面成 $20^\circ$，入射角就是 $70^\circ$。角度数值正确还不够：反射光必须留在入射侧，并沿离开反射点的方向传播。把一条直线的两个相反方向混为一谈，往往会得到“线平行但箭头错误”的光路。

设入射传播方向的单位矢量为 $\vect s$，镜面单位法向为 $\vect n$。分解成沿面与垂直面两部分：
\[
 \vect s=\bigl[\vect s-(\vect s\cdot\vect n)\vect n\bigr]
       +(\vect s\cdot\vect n)\vect n.
\]
反射后沿面分量不变，法向分量反号，因而
\begin{equation}\label{ch27:eq:depth-reflection-vector}
 \vect s_{\rm r}=\vect s-2(\vect s\cdot\vect n)\vect n.
\end{equation}
换用 $-\vect n$ 不改变结果。这一式既可检查平面镜，也可检查曲面镜，只需把 $\vect n$ 换成入射点处的局部法线。球面法线必须经过球心；它一般不是图上的竖直线或水平线。

连续两次反射时，应把第一次反射后的传播方向作为第二次的入射方向。用镜像法展开折线只是帮助计算：镜后虚像和虚线延长线并不表示光实际穿进镜后。求光斑运动时还要先确定最后一条反射线与屏的交点，再对交点位置求时间导数，不能把镜子的角速度直接当作光斑的线速度。
%</block>
%<block:折射定律与多界面光路>
\subsection{多次折射中什么可以消去}
折射角总是相对于当前界面的法线定义。若各界面互相平行，法线方向相同，一段介质内的出射角就等于下一界面的入射角，于是
\[
 n_1\sin i_1=n_2\sin i_2=\cdots=n_k\sin i_k.
\]
中间折射率可以从“最终方向”的关系中消去，但中间层的厚度、侧向位移和光程不能消去。若两界面不平行，例如棱镜的两个侧面，必须先用几何关系换算第二次入射角，再用折射定律。

以厚度为 $t$ 的平行玻璃板为例，板外同为空气，空气侧角为 $i$，板内角为 $r$。第二次折射满足 $n\sin r=\sin i'$，因此 $i'=i$：出射光与入射光同向平行。板内光线的几何长度为 $t/\cos r$，相对原直线的垂直侧移量则为
\[
 b=\frac{t\sin(i-r)}{\cos r}.
\]
这里 $t\tan i-t\tan r$ 是在平行于板面的方向上比较得到的交点差，再乘 $\cos i$ 才得到两条平行光线的最短距离 $b$。当 $n\to1$ 或 $t\to0$ 时，$b\to0$；正入射时也没有侧移，但仍会产生光程变化。

处理球壳或圆柱时，先连“球心或圆心—入射点”得到法线。光线到中心的垂距、光线与半径的夹角和球半径之间满足直角三角形关系；先建立这一关系，才能把折射定律与“碰到内壁”“与内球相切”等条件连接起来。
%</block>
%<block:棱镜偏向、色散与彩虹>
\subsection{最小偏向与色散测量的条件}
棱镜公式中的 $A$ 是两个工作面的夹角，不一定等于图上最醒目的角。设内部两次折射角分别为 $r_1,r_2$，有 $r_1+r_2=A$。在可逆、对称的最小偏向光路中，$r_1=r_2=A/2$、$i_1=i_2=i$，故
\[
 \delta_{\min}=2i-A,\qquad
 n=\frac{\sin[(A+\delta_{\min})/2]}{\sin(A/2)}
\]
（外界为空气）。实验上固定棱镜并只改变波长时，不同颜色一般不同时满足最小偏向条件；不能把同一个入射角问题误做成“每种颜色都重新调整到最小偏向”。

薄棱镜近似还要求 $A,i,r$ 都足够小。把正弦近似成角度后，$i_1\simeq nr_1$、$i_2\simeq nr_2$，于是
\[
 \delta=i_1+i_2-A\simeq(n-1)A,
 \qquad \delta_{\rm violet}-\delta_{\rm red}
 \simeq(n_{\rm violet}-n_{\rm red})A.
\]
这说明色散不仅来自“有棱镜”，更来自不同频率的折射率不同。若需要较大的偏向角或高精度结果，必须回到两次折射的精确关系，不能继续用线性近似。
%</block>
%<block:全反射、临界条件与光纤>
\subsection{从界面临界角推到光纤接收角}
设直光纤芯层折射率为 $n_1$、包层为 $n_2<n_1$，端面外介质为 $n_0$。只讨论穿过光纤轴的子午光线。光从端面进入后与轴成角 $\alpha$；由于端面法线沿轴，端面折射满足
\[
 n_0\sin\theta=n_1\sin\alpha.
\]
在圆柱侧壁处，法线垂直于轴，所以侧壁入射角为 $i=\pi/2-\alpha$。全反射要求
\[
 \sin i=\cos\alpha>\frac{n_2}{n_1}.
\]
将 $\sin^2\alpha=1-\cos^2\alpha$ 代回端面关系，得到
\[
 n_0^2\sin^2\theta<n_1^2-n_2^2.
\]
因此 $\sqrt{n_1^2-n_2^2}$ 称为相应子午光线模型的数值孔径。若其与 $n_0$ 的比值超过1，几何接收角的上限只能是 $90^\circ$，不能对大于1的数取反正弦。等号给临界极限，实际稳定导光还需留出余量；光纤弯曲、端面倾斜与三维斜光线会改变上述简单几何。

这也给出物理检验：$n_1\to n_2$ 时，接收范围缩为零；增大芯层和包层的折射率差，接收范围变大。全反射是界面上不存在传播型折射波，不是说低折射率一侧的电磁场严格为零；后面的隐失场正用来补充这一点。
%</block>
%<block:费马原理与非均匀介质中的光路>
\subsection{驻值方程如何成为折射定律}
将平界面取为横轴，起点在 $(0,a)$、终点在 $(L,-b)$，两侧折射率为 $n_1,n_2$，其中 $a,b>0$。经过界面点 $(x,0)$ 的候选折线路径，其光程为
\[
 \mathcal L(x)=n_1\sqrt{a^2+x^2}+n_2\sqrt{b^2+(L-x)^2}.
\]
求导时，第二段中的 $L-x$ 对 $x$ 求导为 $-1$，因此
\[
 \mathcal L'(x)=n_1\frac{x}{\sqrt{a^2+x^2}}
              -n_2\frac{L-x}{\sqrt{b^2+(L-x)^2}}.
\]
两个分式分别是两侧光线与法线夹角的正弦。令 $\mathcal L'=0$ 就得到斯涅尔定律，而不是一个独立的新折射公式。再求导，
\[
 \mathcal L''(x)=\frac{n_1a^2}{(a^2+x^2)^{3/2}}
                +\frac{n_2b^2}{[b^2+(L-x)^2]^{3/2}}>0,
\]
所以这一具体问题的驻值确实是唯一极小值。一般光学系统中的驻值不一定是最短路；不能把这里的正二阶导数结论推广到所有反射面、焦点和受约束光路。

分层介质中，每一薄层界面均平行，故相对层面法线的角 $i$ 满足 $n\sin i=\text{常量}$。如果改用相对层面的角 $\alpha$，同一规律写成 $n\cos\alpha=\text{常量}$。两式并不矛盾，差别只是角度的起算方向。建立微分方程前，应把图中的角与公式中的角核对一次。
%</block>
